% LaTeX Curriculum Vitae -- WEB VERSION
% Identical to CV_Seo-Won_Chang.tex except that the References section
% lists no names or contact details. Use the full version for applications.
% Seo-Won Chang
% Structure adapted from Jason Blevins' CV template (http://jblevins.org/projects/cv-template/)

\documentclass[letterpaper,11pt]{article}

\usepackage{xcolor}
\usepackage[colorlinks=true,urlcolor=purple,linkcolor=purple!87!black,pdfborder={0 0 0}]{hyperref}
\usepackage{geometry}
\usepackage{etaremune}
\usepackage{enumitem}
\usepackage[T1]{fontenc}
\usepackage[sc,osf]{mathpazo}

\renewcommand{\labelitemi}{$\bullet$}
\renewcommand{\labelitemii}{$\cdot$}

% Set your name here
\def\name{Seo-Won Chang}

\hypersetup{
  colorlinks = true,
  urlcolor = purple,
  pdfauthor = {\name},
  pdfkeywords = {astronomy, astrophysics, time-domain, multi-messenger},
  pdftitle = {\name: Curriculum Vitae},
  pdfsubject = {Curriculum Vitae},
  pdfpagemode = UseNone
}

\geometry{
  body={6.5in, 9.0in},
  left=1.0in,
  top=1.0in
}

\newenvironment{cvitemize}{
  \begin{list}{$\bullet$}{
    \setlength{\leftmargin}{1.5em}
    \setlength{\itemsep}{0.2em}
    \setlength{\parsep}{0em}
  }
}{
  \end{list}
}


\pagestyle{myheadings}
\markright{\name}
\thispagestyle{empty}

\usepackage{sectsty}
\sectionfont{\rmfamily\mdseries\Large}
\subsectionfont{\rmfamily\mdseries\itshape\large}

\setlength\parindent{0em}

% Allow slight line stretching to avoid overfull boxes from URLs / long designations
\setlength{\emergencystretch}{2.5em}
\hyphenation{Django Post-greSQL data-management}

% Tight itemize without bullets for list-style sections
\newenvironment{cvlist}{
  \begin{list}{}{
    \setlength{\leftmargin}{1.5em}
    \setlength{\itemsep}{0.15em}
    \setlength{\parsep}{0em}
  }
}{
  \end{list}
}

% Keep the original bullet-less itemize behaviour
\renewenvironment{itemize}{
  \begin{list}{}{
    \setlength{\leftmargin}{1.5em}
    \setlength{\itemsep}{0.2em}
    \setlength{\parsep}{0em}
  }
}{
  \end{list}
}

\begin{document}

%========================= HEADER =========================
{\huge \name}

\vspace{0.2in}

\begin{minipage}{0.52\linewidth}
  Research Associate Professor \\
  SNU Astronomy Research Center \\
  Seoul National University \\
  Seoul, Republic of Korea
\end{minipage}
\begin{minipage}{0.46\linewidth}
  \begin{tabular}{ll}
    Email: & \href{mailto:seowon.chang@snu.ac.kr}{\tt seowon.chang@snu.ac.kr} \\
    ORCID: & \href{https://orcid.org/0000-0002-3118-8275}{\tt 0000-0002-3118-8275} \\
  \end{tabular}
\end{minipage}

\vspace{0.15in}

%========================= RESEARCH INTERESTS =========================
\section*{Research Interests}

I am an observational astronomer working in \textbf{time-domain and multi-messenger
astronomy}. To study how the sky changes on timescales from minutes to years, I have
spent the past decade building end-to-end infrastructure for wide-field surveys ---
high-precision photometric pipelines, spatio-temporal databases, and real-time
transient-detection systems --- for facilities including SkyMapper, KMTNet, and the
7-Dimensional Telescope (7DT). On these platforms I pursue four connected lines of
research: \textbf{(i) time-domain survey infrastructure and AI-based classification};
\textbf{(ii) stellar magnetic activity} --- flares, starspots, and the
activity--rotation--age relation; \textbf{(iii) stellar pulsation} --- the formation
channels of Blue Large-Amplitude Pulsators (BLAPs); and \textbf{(iv) multi-messenger
astronomy} --- the optical search for gravitational-wave counterparts (kilonovae) and
other high-energy transients. This work is positioned for the era of the Vera C.\ Rubin
Observatory's Legacy Survey of Space and Time (LSST) and the International
Gravitational-Wave Observatory Network (IGWN).

%========================= EDUCATION =========================
\section*{Education}
\begin{itemize}
  \item \textbf{Ph.D. in Astronomy}, Yonsei University, 2014 \\
        Thesis: \emph{Flare and Starspot-induced Variability of Red Dwarf Stars in the
        Open Cluster M37: Photometric Study on Magnetic Activity} \\
        Advisor: Prof.\ Yong-Ik Byun
  \item \textbf{M.S. in Astronomy}, Yonsei University, 2007 \\
        Thesis: \emph{Verification and Characterization of Optical Transient Candidates}
  \item \textbf{B.S. in Astronomy}, Yonsei University, 2005
\end{itemize}

%========================= EMPLOYMENT =========================
\section*{Employment}
\begin{itemize}
\item \textbf{Research Associate Professor}, Research Institute of Basic Sciences,
      Seoul National University (Dec 2023 -- present)
\item \textbf{Research Assistant Professor}, Research Institute of Basic Sciences,
      Seoul National University (Dec 2020 -- Nov 2023)
\item \textbf{OzGrav Postdoctoral Researcher}, Centre for Gravitational Astrophysics,
      Australian National University (Jan 2019 -- Nov 2020)
\item \textbf{CAASTRO Postdoctoral Researcher}, Australian National University
      (Jan 2017 -- Jan 2019)
\item \textbf{Postdoctoral Researcher}, Yonsei University Observatory (Sep 2014 -- Dec 2016)
\item \textbf{Research Assistant}, Near-Earth Space Survey Group, Korea Astronomy and
      Space Science Institute (Jan 2008 -- Dec 2008)
\end{itemize}

%========================= APPOINTMENTS =========================
\section*{Professional Appointments}
\begin{itemize}
\item \textbf{Visiting Fellow}, Research School of Astronomy \& Astrophysics,
      Australian National University (Nov 2020 -- Nov 2025; Feb 2026 -- Feb 2027)
\end{itemize}

%========================= GRANTS =========================
\section*{Research Grants as Principal Investigator}
\begin{itemize}
\item \textbf{The Origins of Blue Large-Amplitude Pulsators (BLAPs): Binary Interaction,
      Merger, and Supernova Survivor Scenarios}, National Research Foundation of Korea
      (Core Research Program), Mar 2026 -- Feb 2029 
\item \textbf{Understanding Galactic Disc Evolution with Stellar Activity and Kinematic
      Age}, National Research Foundation of Korea, Jun 2023 -- May 2026 
\end{itemize}

%========================= AWARDS =========================
\section*{Honors \& Awards}
\begin{itemize}
\item \textbf{Top 100 National R\&D Achievements of 2025} (team award),
      Ministry of Science and ICT, Republic of Korea, January 2026 \\
      ``Commencement of Observations with the 7-Dimensional Telescope (7DT),
      the World's First Wide-Field Simultaneous Spectral-Imaging Multiple
      Telescope System'' \\
      \textit{Team:} Myungshin Im (PI), Hyung Mok Lee, Ji Hoon Kim,
      \textbf{Seo-Won Chang}, Gregory S. H. Paek, and Hyunho Choi
\item \textbf{Yonsei Thesis Award} for Excellence in Graduate Research,
      Yonsei University, 2014
%\item \textbf{Yonsei Alumni Fellowship}, Yonsei University, 2004
\end{itemize}

%========================= RESEARCH HIGHLIGHTS =========================
\section*{Research Highlights}
\emph{Four connected research lines, emphasizing recent and ongoing directions (since 2022).}

\subsection*{Time-domain Astronomy: Survey Infrastructure and AI-based Classification}
I build the infrastructure that enables time-domain discovery --- high-precision photometry
pipelines, spatio-temporal databases, and the public release of wide-field surveys, most
recently as first-author lead of the KS4 Data Release~1 \emph{(Chang et al.\ 2026, JKAS)}.

\smallskip
On the methodological side, I am increasingly focused on machine-learning classification
--- variable types, periodic variables, and real/bogus transients \emph{(Lee et al.\ 2025,
AJ)} --- with the goal of decision-aware AI for LSST-scale time-domain astronomy: systems
that not only classify but help decide, in real time, which of millions of nightly alerts
merit follow-up.

\smallskip
In parallel, I am pursuing time-domain studies based on single- and multi-epoch SEDs,
combining 7DS medium-band spectrophotometry with SPHEREx infrared spectral surveys to
classify and characterize variable and transient sources across the optical--infrared
range \emph{(Paek et al.\ 2026, ApJ; Kim et al.\ 2026, AJ)}.

\subsection*{Stellar Magnetic Activity: Flares, Starspots, and the Activity--Age Relation}
My work in this area began with magnetically active M dwarfs --- their flares and
starspot-induced rotational variability --- and the age--rotation--activity relation,
established through large-scale photometric measurements including the M-dwarf flaring
fraction and a spatially dependent activity--kinematic-age relation across the solar
cylinder \emph{(Chang et al.\ 2020, 2022, MNRAS; first author)}.

\smallskip
I am now extending this work to a statistical study of flare activity in FGK dwarfs
through my NRF-funded program, and further toward spectroscopic detection of coronal mass
ejections and stellar storms --- the space-weather events that shape the environments of
exoplanets around Sun-like and other host stars \emph{(Oh et al., JKAS, in revision;
corresponding author)}.

\smallskip
Looking forward, I am particularly interested in simultaneous multi-wavelength studies of
flares --- combining optical photometry, emission/absorption-line spectroscopy, X-ray, and
radio observations. The Einstein Probe opens new opportunities here: many of its X-ray
transient alerts are in fact stellar flares, well suited to joint X-ray--optical studies
with optical time-domain facilities.

\subsection*{Stellar Pulsation: Blue Large-Amplitude Pulsators}
Blue Large-Amplitude Pulsators (BLAPs) are a recently recognized class of hot pulsating
stars --- nearly 200 have been found toward the Galactic bulge since their 2017 discovery
by OGLE. Their evolution requires the loss of most of the stellar envelope, yet unlike hot
subdwarfs (sdBs), which are predominantly found in binaries, BLAPs show few binary
companions. their formation channel remains actively debated. %, and no homogeneous spectroscopic characterization of the class yet exists.

\smallskip
I discovered a new BLAP in SkyMapper and characterized it as a helium-deficient,
thick-disk star --- a candidate for an alternative formation channel
\emph{(Chang et al.\ 2024, MNRAS; first author)}. This discovery anchors my NRF Core
Research program (2026--2029) on BLAP formation channels. Combining Gemini spectroscopy
with stellar-atmosphere modeling, the program aims to build the homogeneous spectroscopic
sample the field lacks and to map the connection between BLAPs and sdB/sdBV pulsators.

\smallskip
More broadly, this work extends toward the study of post-mass-transfer binaries ---
the stripped-envelope systems that link BLAPs, hot subdwarfs, and the broader landscape
of binary stellar evolution.

\subsection*{Multi-Messenger Astronomy: Optical Counterparts of Gravitational Waves}
From early 2017 through the O4 run, I devoted myself to the optical search for
gravitational-wave counterparts --- participating in the GW170817/AT2017gfo discovery,
building and operating the SkyMapper alert follow-up pipeline through LIGO/Virgo O3
\emph{(Chang et al.\ 2021, PASA; first author)}, and continuing follow-up with KMTNet and
7DT within GECKO.

\smallskip
My focus is now expanding from follow-up operations to the science of the sources
themselves: kilonovae from NS--NS and NS--BH mergers and from gamma-ray bursts
\emph{(Yang et al.\ 2024, Nature; Pillas et al.\ 2025, Phys.\ Rev.\ D)}, single-epoch
SED-based identification of orphan kilonovae in the Rubin/LSST alert stream without
gravitational-wave triggers \emph{(Paek et al.\ 2026, ApJ)}, and their connection to
host-galaxy environments.

\smallskip
Looking ahead to the IGWN (IR1/O5) era, I am also interested in developing real-time
automated spectroscopic follow-up systems and $\gamma$-ray CubeSat instrumentation for
multi-messenger astronomy.


%========================= INFRASTRUCTURE & SOFTWARE =========================
\section*{Infrastructure, Software \& Technical Skills}
\begin{cvitemize}
\item \textbf{Survey data release \& systems:} End-to-end development and public release of
      wide-field survey data --- architected the KS4 DR1 database (catalog merging,
      quality-assurance flagging, HiPS imagery) and delivered it to the international
      community via NOIRLab Astro Data Lab and CDS/Aladin \& VizieR; lead developer of the 7DT
      data-management portal (full-stack Django, PostgreSQL/PostGIS).
\item \textbf{Time-series \& transient pipelines:} High-precision multi-aperture
      photometry, noise de-trending, variability detection, differential image analysis,
      and multiband period analysis.
\item \textbf{Machine learning:} Classification problems --- variable-star types, periodic
      variables, and real/bogus transients; outlier detection in spectrophotometric
      multi-epoch SEDs; and emulators for stellar-atmosphere models.
\item \textbf{Tools:} Python (primary), full-stack web development, and Virtual Observatory
      tools
\end{cvitemize}

%========================= SURVEYS & COLLABORATIONS =========================
\section*{Survey \& Collaboration Participation}
\begin{cvitemize}
\item \textbf{7-Dimensional Telescope (7DT) / 7-Dimensional Sky Survey (7DS)} ---
core member; data portal and database development; multi-messenger follow-up; time-domain
science
\item \textbf{KMTNet} --- KS4 Data Release 1 lead; transient and neutrino ToO programs
\item \textbf{SkyMapper Southern Survey} --- variability and transient science (DR1--DR4);
multi-wavelength stellar-activity collaboration (GALAH, eROSITA, ASKAP)
\item \textbf{Korean Gravitational Wave Group} --- gravitational-wave optical follow-up within GECKO (GW-EM Counterpart Korean Observatory)
\item \textbf{OzGrav} (ARC Centre of Excellence for Gravitational Wave Discovery) ---
SkyMapper-side lead for gravitational-wave optical follow-up
\item \textbf{SPHEREx (Korean Legacy Science Program)} --- proposed two legacy science
      projects: infrared activity patterns in late-M to early-L dwarfs (time-domain study
      with SPHEREx and SkyMapper), and the search for wide-orbit companions of Blue
      Large-Amplitude Pulsators through infrared excess
\item \textbf{Rubin/LSST}  --- member, Transients and Variable Stars (TVS) and
      Stars, Milky Way and Local Volume (SMWLV) Science Collaborations; Junior Associate,
      Korean Rubin Science Collaboration; member of the FINK community alert broker
\end{cvitemize}

%========================= TELESCOPE TIME =========================
\section*{Telescope Time (as Principal Investigator)}
\begin{cvitemize}
\item Gemini GPP XT1 Program, 2025B, G-2025B-0498-V (Band-2, 5.2\,hr)
\item ToO Spectroscopy of Gravitational-Wave Optical Counterparts with SkyMapper,
2018Q4--2019Q3
\item SkyMapper Transient Follow-up Program, 2017B--2019A
\end{cvitemize}
\emph{Co-Investigator on KMTNet time-domain and ToO follow-up programs (gravitational-wave,
neutrino, and transient alerts; 30-minute time-domain survey; 2020--2029) and Gemini
GW follow-up (2023--2026).}

%========================= PUBLICATIONS =========================
\section*{Publications}
First-author or corresponding-author papers are marked with $\star$. Author role is
indicated after each entry. Full and up-to-date metrics are available via
\href{https://orcid.org/0000-0002-3118-8275}{ORCID} and ADS.

\subsection*{Refereed Publications}
{\sloppy
\begin{etaremune}
\item \emph{Masking Algorithm for CCD Bleeding in Korea Microlensing Telescope Network
Images}, Shin, J., Jeong, M., Im, M., \textbf{Chang, S.-W.} et al.\ 2026h, JKAS,

\item $\star$ \emph{High-Resolution Spectroscopic Analysis of Chromospheric Line Evolution
during an Energetic Flare on AD~Leo}, Oh et al.\ 2026g, JKAS, [corresponding author]

\item \emph{7DT Insight: Variability in Young Stellar Objects}, Kim, M.-R., Lee, J.-E.,
Im, M., \textbf{Chang, S.-W.} et al.\ 2026f, AJ [co-author] %\href{https://arxiv.org/abs/2605.18366}{arXiv:2605.18366}

\item $\star$ \emph{KMTNet Synoptic Survey of Southern Sky III: The First Data Release},
\textbf{Chang, S.-W.} et al.\ 2026e, JKAS, 59, 1
[first author] 
\item \emph{A Hybrid Framework for Kilonova Anomaly Detection Using Single-epoch SEDs from the 7-Dimensional Telescope}, Paek et al.\ 2026d, ApJ
[co-author] 

\item \emph{Pre-perihelion Emergence of the CN Gas Coma in 3I/ATLAS Resolved by the
7-Dimensional Telescope}, Paek et al.\ 2026c, ApJ [co-author] 

\item \emph{KMTNet Synoptic Survey of Southern Sky II: Data Reduction Pipeline for
Reference Imaging and Real-time Transient Detection}, Jeong et al.\ 2026b, JKAS
[co-author]

\item \emph{The Redshifts from 122 Bands: Comparative Redshift Forecast for Low-Resolution
Spectra from SPHEREx and 7DS}, Bae et al.\ 2026a, A\&A [co-author] 

\item \emph{Investigating the Effects of Point Source Injection Strategies on
KMTNet Real/Bogus Classification}, Lee, D., Paek, G.~S.~H., \textbf{Chang, S.-W.} et al.\
2025, AJ [co-author; key contributor] 

\item \emph{AllBRICQS: The Discovery of Luminous Quasars in the Northern Hemisphere},
Choi et al.\ 2025, ApJS [co-author]

\item \emph{Limits on the Ejecta Mass and Binary Properties During the Search for Kilonovae
Associated with Neutron Star--Black Hole Mergers: S230518h, GW230529, S230627c, and
S240422ed}, Pillas et al.\ 2025, Phys.\ Rev.\ D [multi-messenger collaboration]

\item \emph{GECKO Follow-up Observation of the Binary Neutron Star--Black Hole Merger
Candidate S230518h}, Paek et al.\ 2025, ApJ [co-author]

\item \emph{Exploring Unobscured Quasi-stellar Objects in the Southern Hemisphere with
KS4}, Kim, Y.~J. et al.\ 2024, ApJS, 275, 46 [co-author]

\item \emph{SkyMapper Southern Survey: Data Release 4}, Onken et al.\ 2024, PASP, 41, 61
[co-author]

\item $\star$ \emph{Discovery of a New Blue Large-Amplitude Pulsator in SkyMapper DR2:
SMSS J184506.82$-$300804.7}, \textbf{Chang, S.-W.} et al.\ 2024, MNRAS
[first author] 

\item \emph{A Lanthanide-rich Kilonova in the Aftermath of a Long Gamma-ray Burst},
Yang et al.\ 2024, Nature [multi-messenger collaboration] 

\item $\star$ \emph{Spatially Dependent Photometric Activity of M Dwarfs in the Solar
Cylinder}, \textbf{Chang, S.-W.} et al.\ 2022, MNRAS, 517, 2842 [first author]

\item \emph{Multi-filter Photometry of Solar System Objects from the SkyMapper Southern
Survey}, Sergeyev et al.\ 2022, A\&A [co-author] 

\item \emph{GECKO Optical Follow-up of Three Binary Black Hole Merger Events
(GW190408\_181802, GW190412, GW190503\_185404)}, Kim et al.\ 2021, ApJ [co-author]

\item $\star$ \emph{SkyMapper Optical Follow-up of Gravitational Wave Triggers: Alert
Science Data Pipeline and LIGO/Virgo O3 Run}, \textbf{Chang, S.-W.} et al.\ 2021, PASA
[first author]

\item \emph{A Spectroscopically Confirmed Gaia-selected Sample of 318 New Young Stars
within $\sim$200~pc}, Zerjal et al.\ 2021, MNRAS [co-author; observations]

\item \emph{Neutron Star Extreme Matter Observatory: A Kilohertz-band Gravitational-wave
Detector in the Global Network}, Ackley et al.\ 2020, PASA, 37, 47
[OzGrav collaboration]

\item \emph{Probing the Extragalactic Fast Transient Sky at Minute Timescales with DECam},
Andreoni et al.\ 2020, MNRAS, 491, 5852 [co-author; flare light-curve interpretation]

\item $\star$ \emph{Photometric Flaring Fraction of M Dwarf Stars from the SkyMapper
Southern Survey}, \textbf{Chang, S.-W.}, Wolf, C.\ \& Onken, C.~A.\ 2020, MNRAS, 491, 39
[first author]

\item \emph{SkyMapper Southern Survey: Second Data Release (DR2)}, Onken et al.\ 2019,
PASA, 36, e033 [co-author; noise filtering, quality cuts, variability index]

\item \emph{Five New Real-time Detections of Fast Radio Bursts with UTMOST},
Farah et al.\ 2019, MNRAS, 488, 2989 [co-author; optical follow-up]

\item \emph{A Fast Radio Burst with Frequency-dependent Polarization Detected during
Breakthrough Listen Observations}, Price et al.\ 2019, MNRAS, 486, 3636
[co-author; optical follow-up]

\item $\star$ \emph{New Photometric Pipeline to Explore Temporal and Spatial Variability
with KMTNet DEEP-South Observations}, \textbf{Chang, S.-W.} et al.\ 2018, JKAS, 51, 129
[first author]

\item \emph{Detecting Variability in Massive Astronomical Time-series Data III: Variable
Candidates in SuperWASP DR1}, Shin, M.-S., \textbf{Chang, S.-W.} et al.\ 2018, AJ, 156, 5
[co-author]

\item \emph{A Nearby Superluminous Supernova with a Long Pre-maximum Plateau and Strong
C~II Features}, Anderson et al.\ 2018, A\&A, 620, 67 [co-author; SkyMapper data]

\item \emph{FRB Microstructure Revealed by the Real-time Detection of FRB170827},
Farah et al.\ 2018, MNRAS, 478, 1209 [co-author; optical follow-up]

\item \emph{SkyMapper Southern Survey: First Data Release (DR1)}, Wolf et al.\ 2018,
PASA, 35, 10 [co-author; quality cuts, time-domain variability]

\item \emph{Follow-up of GW170817 and its Electromagnetic Counterpart by Australian-led
Observing Programs}, Andreoni et al.\ 2017, PASA, 34, 21
[multi-messenger collaboration; SkyMapper follow-up]

\item \emph{Multi-messenger Observations of a Binary Neutron Star Merger},
Abbott et al.\ 2017, ApJL, 848, 2 [multi-messenger collaboration; SkyMapper follow-up]

\item $\star$ \emph{Photometric Study on Stellar Magnetic Activity I: Flare Variability of
Red Dwarf Stars}, \textbf{Chang, S.-W.}, Byun, Y.-I.\ \& Hartman, J.~D.\ 2015, ApJ, 814, 35
[first author]

\item $\star$ \emph{A New Catalog of Variable Stars in the Field of the Open Cluster M37},
\textbf{Chang, S.-W.}, Byun, Y.-I.\ \& Hartman, J.~D.\ 2015, AJ, 150, 27 [first author]

\item $\star$ \emph{A New Method for Robust High-Precision Time-Series Photometry from
Well-Sampled Images: Archival MMT/Megacam Observations of M37}, \textbf{Chang, S.-W.},
Byun, Y.-I.\ \& Hartman, J.~D.\ 2015, AJ, 149, 135 [first author]

\item \emph{The EPOCH Project I: Periodic Variable Stars in the EROS-2 LMC Database},
Kim et al.\ 2014, A\&A, 566, 43 [co-author; data and database management]

\item $\star$ \emph{Statistical Properties of Galactic $\delta$ Scuti Stars: Revisited},
\textbf{Chang, S.-W.} et al.\ 2013, AJ, 145, 132 [first author]

\item \emph{Detecting Variability in Massive Astronomical Time-series Data II: Variable
Candidates in the Northern Sky Variability Survey}, Shin et al.\ 2012, AJ, 143, 65
[co-author]

\item \emph{BVR Photometry of Supergiant Stars in Holmberg~II}, Sohn, Y.-J.,
\textbf{Chang, S.-W.} et al.\ 2006, JASS, 23, 1 [co-author; data analysis]

\end{etaremune}
\par}

\subsection*{Submitted / In Revision}
\begin{etaremune}
\item \emph{KS4GC: A Red-Sequence-Based Galaxy Cluster Catalog in the Southern Sky from
KS4 DR1}, Park et al.\ 2026, A\&A, submitted [co-author]
\item \emph{Construction of spatially resolved star formation and 3.3~$\mu$m PAH distribution map using the combined 7DT and SPHEREx photometric spectra}, Shim et al.\ 2026, ApJ, submitted
\item \emph{Optical Transmission Spectroscopy of the Hot Jupiter WASP-74b Using Medium-Band Filters with the 7-Dimensional Telescope}, Bae et al.\ 2026, A\&A, submitted

\end{etaremune}

\subsection*{Selected Conference Proceedings}
\begin{etaremune}
% Py7DT: Data Reduction Pipeline of the 7-Dimensional Telescope, 2026, SPIE
\item \emph{The 7-Dimensional Telescope and the 7-Dimensional Sky Survey: Status Report on Final Commissioning and Survey Operation}, Kim, J.~H., Im, M., Lee, H.~M.\ \& \textbf{Chang, S.-W.} et al. 2026, SPIE
\item \emph{Optical Analysis of SN 2023ixf: From 10 Years Before to 100 Days After the
Explosion}, Kim, S., Im, M., Paek, G.~S.~H., Lim, G., \textbf{Chang, S.-W.} et al.\ 2024,
The Twelfth Pacific Rim Conference on Stellar Astrophysics, ASP Conf.\ Ser.\ 536, 61
\item \emph{Project Management for Ground-based Telescope Array Development},
Kim, J.~H., Im, M., Lee, H.~M.\ \& \textbf{Chang, S.-W.} 2024, SPIE
\item \emph{Introduction to the 7-Dimensional Telescope: Commissioning Procedures and Data
Characteristics}, Kim, J.~H. et al.\ 2024, SPIE 13094
\item \emph{Spatially Dependent Photometric Activity of M Dwarfs in the Solar Cylinder},
\textbf{Chang, S.-W.}, Wolf, C.\ \& Onken, C., 2022, Cool Stars 21 (CS21)
\end{etaremune}

\subsection*{GCN Circulars \& Astronomical Telegrams}
Lead- or co-author on numerous GCN Circulars and Astronomer's Telegrams (ATels) reporting
real-time optical follow-up of gravitational-wave events, gamma-ray bursts, and other
transients (e.g., fast radio bursts), through the SkyMapper, ANU 2.3\,m
WiFeS, KMTNet/GECKO, and 7DT programs. Full list available via ADS.

%========================= TALKS =========================
\section*{Selected Talks \& Presentations}
\emph{Representative presentations on my recent (last five years) research topics.}
\begin{etaremune}
\item \emph{Searching for Blue Large-Amplitude Pulsators: SkyMapper Results and
Strategies}, 3rd Stars, Milky Way, Local Volume and Exoplanet (SMiLE) Workshop, 2026
\item \emph{KMTNet Synoptic Survey of Southern Sky: The First Data Release (KS4 DR1)},
Korean Astronomical Society Meeting, 2026
\item Session Chair, Special Session \emph{Early Science Results from 7DT}, Korean
Astronomical Society Spring Meeting, 2026
\item \emph{Searching for Stellar Flares from FGKM Dwarfs in SkyMapper DR4}, Korean
Astronomical Society Meeting, 2025
\item \emph{SkyMapper's Fourth Data Release and Stellar Science Cases}, Stars, Milky Way,
Local Volume and Exoplanet (SMiLE) Workshop, 2024
\item \emph{KMTNet Synoptic Survey of Southern Sky: Early Data Release Preparation},
Korean Astronomical Society Meeting, 2023
\item \emph{Evidence for Spatially Dependent Magnetic Activity of M Dwarfs in the Solar
Cylinder}, Korean Astronomical Society Meeting, 2022
\item \emph{Preparation for the Advanced LIGO-Virgo-KAGRA O4 Observing Run}, Centre for
Gravitational Astrophysics / RSAA Joint Seminar Series, ANU, 2021 (colloquium)
\item \emph{SkyMapper Optical Follow-up of Gravitational Wave Triggers: Alert Science
Data Pipeline and LIGO/Virgo O3}, European Astronomical Society Annual Meeting, 2021
\item \emph{SkyMapper Optical Follow-up of Gravitational Wave Triggers and OzGrav~2.0},
Seoul National University, 2021 (colloquium)

\end{etaremune}


%========================= TEACHING & MENTORING =========================
\section*{Teaching \& Mentoring}
\begin{cvitemize}
\item Thesis Committee Member, M.S. thesis on time-resolved spectroscopy of an
      energetic flare on AD~Leo, Seoul National University, 2025
\item Mentor, high-school research project (\emph{Efficient Linkage Algorithms for
      DEEP-South Asteroid Discovery}), Hansung Science High School, 2022
\item Advisor, Undergraduate Summer/Winter Research Programs, Yonsei University
      (2012, 2013, 2015)
\item Teaching Assistant, Yonsei University: \emph{Galaxies and the Universe} (2006, 2007),
      \emph{Astrophysics II} (2005), \emph{Understanding of Space} (2005)
\end{cvitemize}

%========================= SERVICE =========================
\section*{Professional Service}
\begin{cvitemize}
\item Member, Multi-Messenger Astronomy Planning Task Force, Korea Astronomy and Space
      Science Institute (KASI) --- contributed to a national research-planning proposal,
      \emph{Extreme-Gravity Astrophysics and Cosmic Expansion through Gravitational-Wave
      and Multi-Messenger Observations}
\item Committee Member, Working Group on Multi-Messenger Observing Facilities and
      Data/Computing Infrastructure, \emph{The Decadal Survey for Korean Astronomy
      (2026--2035)}
\item Referee/Reviewer: AAS Journals (ApJ, AJ, ApJS), Journal of the Korean Astronomical
      Society (JKAS), Publications of the Korean Astronomical Society (PKAS)
\end{cvitemize}


%========================= OUTREACH =========================
\section*{Outreach \& Public Engagement}
\begin{cvitemize}
\item Led the public release of the KS4 DR1 dataset to the global astronomical community,
      announced through coordinated releases by KASI, SNU,
      \href{https://cds.unistra.fr/blog/2026/04/29-ks4-data-release/}{CDS}, and
      \href{https://noirlab.edu/science/news/announcements/sci26025}{NOIRLab} (2026)
\item Contributing Author, \emph{Encyclopedia of Astronomy}, published by the Korean
      Astronomical Society via Naver Knowledge Encyclopedia
\end{cvitemize}

%========================= RESEARCH PROJECT PARTICIPATION =========================
%\section*{Selected Research Project Participation}
%\begin{itemize}
%\item Multi-Messenger Astronomy with the 7-Dimensional Telescope and Artificial
%     Intelligence (Jul 2021 -- Dec 2025)
%\item Multi-Messenger Astronomy and Massive Astrophysical Systems through Wide-area
%     Imaging Survey (Jun 2020 -- Feb 2025)
%\end{itemize}

%========================= REFERENCES =========================
\section*{References}
Available on request.

\end{document}
